\chapter{Stability and Control Characteristics}
\label{ch:scc}

\textcolor{red}{\cite{NASA_SP8007_2019} used for critical buckling loads: }

\todo[inline]{From the \textit{Project Guide}: \\ Stability and control characteristics include e.g. control forces, CoG limits, attitude control torques, momentum storage and unload capacity, disturbance torques, control frequency, etc.\\
A first version of design driving characteristics is included in the Mid-Term Report, a final, extended version including all relevant characteristics is given in the Final Report.}

\todo[inline]{Chapter introduction pending. Remember to connect it to the previous chapters.}


\section{Control Forces}
\label{sec:control_forces}
The control forces on each of the control surfaces were calculated using the Flow5 software. This software uses a potential flow solver, viscosity and wake models for preliminary estimation of aircraft characteristics. These include lift distribution, drag estimation, static stability and dynamic stability analysis. \footnote{\href{https://flow5.tech/docs/flow5_doc/flow5_doc.html}{Flow5 methods overview}(Accessed on 08.06.2026)} The aircraft was modeled meshed in this software using the fuselage CAD model and a built-in wing generation tool. Control surface size was estimated based on the reference aircraft. The output from Flow5 will give a range of moment values for the actuators needed for the control surface deflection, which are essential for selection actuators for these subsystems. Furthermore, the Flow5 aircraft model is used together with its inertial properties coming from the CAD model to estimate dynamic stability envelope.
\begin{figure}[H]
    \centering
    \includegraphics[width=0.8\linewidth]{}
    \caption{HUGO aircraft model for analysis in Flow5}
    \label{fig:flow5_overview}
\end{figure}
\subsection{Flight conditions}
\autoref{} shows the control surface characteristics of the reference aircraft.
\begin{table}[]
    \centering
    \begin{tabular}{c|c}
         &  \\
         & 
    \end{tabular}
    \caption{Caption}
    \label{tab:ctrl_surfaces_ref_aircraft}
\end{table}
Based on this data, the following control surfaces were investigated in the Flow5 model:
\begin{table}[]
    \centering
    \begin{tabular}{c|c}
         &  \\
         & 
    \end{tabular}
    \caption{Caption}
    \label{tab:hugo_ctrl_surfaces}
\end{table}

These control surfaces were tested in an approach speed regime where $V_{appr} = m \cdot s^{-1}$. At this flight velocity, the hinge moment on the ailerons is ... at the maximum deflection with both ailerons deflected in opposite directions. The resulting roll rate p is calculated at ... The turn radius is ... \\

At the target experiment flight speed at $M = 0.75$, the required aileron deflection to reach the approach roll rate is ... This deflection generates a hinge moment of ... This value was obtained in Flow5 and then scaled using the Prandtl-Glauert compressibility correction
\subsection{Control surface analysis verification}
In order to verify the control surface moments, a simulation setup in identical flight conditions was used.
\subsection{Dynamic characteristics of the aircraft}

\section{Centre of Gravity Limits}
\label{sec:centre_of_gravity_limits}

\todo[inline]{Stability \& Controllability, scissor plot etc.}











\todo[inline]{Chapter conclusion pending. Remember to connect it to the chapters that come after.}